% =============================================================================
% SECTION 2A: RELATED WORK AND NOVELTY
% =============================================================================

\section{Related Work and Contributions}

\subsection{Existing Approaches}

Current construction takeoff solutions fall into three categories:

\textbf{Manual/Semi-Automated Tools.}
Traditional tools like PlanSwift~\cite{planswift2024} and Bluebeam require human operators to trace elements and apply counts.
While accurate, they scale poorly and depend on operator expertise.

\textbf{Rule-Based Estimation.}
Commercial systems (e.g., STACK~\cite{stack2024}) embed framing rules in code, limiting adaptability across jurisdictions and practices.
Rule updates require software modifications, and decisions lack auditability.

\textbf{AI-Powered Takeoff.}
Recent AI tools (e.g., Togal.AI~\cite{togalai2024}) use computer vision to detect areas and count symbols~\cite{pizarro2022floorplan_survey, jamieson2025symbol_detection},
but typically output quantities without assembly reasoning, cut optimization, or code compliance verification.

\subsection{Limitations of Prior Work}

\begin{enumerate}
  \item \textbf{Embedded knowledge}: Rules and practices are hard-coded, not inspectable or versionable
  \item \textbf{Missing provenance}: BOM line items cannot be traced to plan facts and applied rules
  \item \textbf{No optimization}: Stock selection and cutting are left to downstream processes
  \item \textbf{Limited compliance}: Code citations require manual verification
  \item \textbf{No learning}: Systems do not improve from project outcomes
\end{enumerate}

\subsection{Key Contributions}

This proposal presents the following novel contributions:

\begin{enumerate}
  \item \textbf{KG-Centered Architecture}: A knowledge graph backbone that externalizes construction rules, material catalogs, and historical precedents---enabling inspection, versioning, and auditing~\cite{pan2024unifying_kg_llm}.

  \item \textbf{Provenance-First Design}: Every BOM line item links to (a) plan facts, (b) applied inference rules, and (c) code references, supporting explainability and regulatory compliance.

  \item \textbf{Multi-Agent Reasoning}: Specialized agents for extraction QA, component inference, compliance checking, and instruction generation~\cite{li2025agentic_survey, zhang2025tool_learning}, with clear separation from deterministic optimization.

  \item \textbf{Integrated Cut Optimization}: Stock selection and cutting as a first-class concern, using ILP solvers~\cite{ortools2024} with kerf handling and multi-stock support~\cite{gilmore1961cutting}.

  \item \textbf{Continuous Improvement Loop}: Human corrections and project outcomes feed back into the KG as new precedents, enabling system-wide learning.
\end{enumerate}

\subsection{Differentiation Summary}

Table~\ref{tab:differentiation} compares our approach to existing solutions.

\begin{table}[ht]
\centering
\caption{Comparison with Existing Approaches}
\label{tab:differentiation}
\begin{tabular}{lcc}
\toprule
\textbf{Capability} & \textbf{Existing} & \textbf{Ours} \\
\midrule
Knowledge storage & Code-embedded & KG-externalized \\
Rule versioning & Manual updates & Temporal validity \\
Provenance & None/limited & Full chain \\
Cut optimization & Separate/none & Integrated \\
Compliance & Manual check & Auto w/ citations \\
Learning & Static & Continuous \\
\bottomrule
\end{tabular}
\end{table}
